% Part: lambda-calculus
% Chapter: lambda-definability
% Section: partial-recursive-functions

\documentclass[../../../include/open-logic-section]{subfiles}

\begin{document}

\olfileid{lam}{ldf}{par}
\olsection{పాక్షిక పునరావృత్త ప్రమేయాలు \usetoken{S}{lambda definable}}

పాక్షిక పునరావృత్త ప్రమేయాలు అంటే ప్రాథమిక ప్రమేయాల నుంచి
సంయుక్తం, ఆదిమ పునరావృత్తి, అపరిమిత కనిష్ఠీకరణ ద్వారా
పొందగల ప్రమేయాలు. అపరిమిత అన్వేషణలో వాడే ప్రమేయాలు
సక్రమమైనవి కావాలని ఇక్కడ కోరరు; ఈ విషయంలో ఇవి సామాన్య
పునరావృత్త ప్రమేయాలకంటే భిన్నం. సక్రమతను కోరనందువల్ల,
కనిష్ఠీకరణ ద్వారా నిర్వచించిన ప్రమేయాలు కొన్ని సందర్భాల్లో
నిర్వచితం కాకపోవచ్చు.

మొదట చూస్తే, సంపూర్ణమైన సామాన్య పునరావృత్త ప్రమేయాలన్నీ
\tetoken{లాంబ్డాతో నిర్వచించదగినవి}{!!{lambda definable}}
అని చూపడానికి వాడిన పద్ధతులనే అన్ని పాక్షిక పునరావృత్త
ప్రమేయాలు \tetoken{లాంబ్డాతో నిర్వచించదగినవి}{!!{lambda definable}}
అని నిరూపించడానికి కూడా వాడవచ్చనిపిస్తుంది. ఉదాహరణకు, $f$, $g$లను
వరుసగా $F$, $G$ పదాలు
\tetoken{లాంబ్డాతో నిర్వచిస్తే}{!!{lambda define}},
$f$తో $g$ సంయుక్తాన్ని $\lambd[x][F (G x)]$
\tetoken{లాంబ్డాతో నిర్వచిస్తుంది}{!!{lambda define}}.
కానీ ప్రమేయాలు పాక్షికమైనప్పుడు చిక్కు వస్తుంది. $g(x)$
నిర్వచితం కాకపోతే, $G x$కు నియత రూపం ఉండదన్న మాట.
సాధారణంగా అప్పుడు $F (G x)$కు కూడా నియత రూపం ఉండదు;
మనకు కావలసిందీ అదే. అయితే $F$ అనేది
$\lambd[x][\lambd[y][y]]$ అయిన సందర్భాన్ని చూడండి.
అప్పుడు దాని ఆర్గ్యుమెంటుకు నియత రూపం లేకపోయినా $F (G x)$కు
$\lambd[y][y]$ అనే నియత రూపం ఉంటుంది.

ఈ చిక్కును అధిగమించవచ్చు. నిర్వచితం కాని విలువలు నియత
రూపం లేని పదాల ద్వారానే సూచించబడేలా అన్ని పాక్షిక పునరావృత్త
ప్రమేయాలను \tetoken{లాంబ్డాతో నిర్వచించే}{!!{lambda define}}
మార్గాలు ఉన్నాయి. అయితే అవి సామాన్య పునరావృత్త ప్రమేయాల
కోసం మనం వాడిన పద్ధతికంటే కొంత సంక్లిష్టమైనవి, తక్కువ
సహజమైనవి. అందుకే ఇక్కడ నిరూపణ లేకుండా సిద్ధాంతాన్ని
నమోదు చేస్తున్నాం:

\begin{thm}
  అన్ని పాక్షిక పునరావృత్త ప్రమేయాలు
  \tetoken{లాంబ్డాతో నిర్వచించదగినవి}{!!{lambda definable}}.
\end{thm}

\end{document}
